\section{Introduction}
\label{sec:introduction}

Graphic design communicates information through the coordinated use of typography, imagery, color, and spatial organization \citep{huang2023survey,zou2025fragment}.
Creating a coherent composition involves decisions about what content to present, how to establish visual hierarchy, and how to adapt a design to its intended medium.
Recent research has brought generative methods to a growing range of such tasks, including posters and advertisements \citep{posterlayout2023,cglgan2022}, presentation slides \citep{autopresent2025,pptagent2025}, infographics \citep{bizgen2025}, and scientific illustrations \citep{autofigure2026}.
Despite their different conventions, these tasks share a need to align semantic content with visual organization.

A central challenge is that a graphic design is not only an image to be viewed, but often an artifact to be revised.
Designers may replace text, reposition visual elements, adjust hierarchy, or adapt a composition to a new format while retaining other parts of the work.
Generative approaches address different parts of this process, from spatial composition \citep{layoutgan2019,layoutdm2023} to fuller design synthesis and structured representations \citep{cole2023,opencole2024,designascode2026}.
Yet the appearance of the rendered result alone does not reveal what can be changed afterward.
Two visually similar designs may have very different underlying structures and, consequently, very different affordances for further editing.

This distinction matters because \emph{editability} is not a single property.
A format may expose individual elements without providing a reliable way to carry out a requested change, while an image-based editing system may satisfy a visual instruction without delivering separable objects.
Some edits also require dependent changes to neighboring content, so preserving everything outside a target region is not always the appropriate objective \citep{yang2026infoedit}.
Assessing revision therefore requires attention to the delivered artifact, the operations it permits, and the evidence that those operations succeed.

Evaluation practices make these differences difficult to compare.
Studies of spatial composition often report geometric or distributional quality \citep{layoutdm2023}, whereas other work assesses perceptual alignment, usability, or instruction-following edits \citep{vectorfusion2023,vinci2021,svgeditbench2024}.
These measurements can be valuable for their intended tasks, but they do not establish the same claim about a design's usability after generation.
In particular, visual fidelity, structural editability, and measured editing performance should not be treated as interchangeable outcomes.

Existing surveys offer broad perspectives on graphic design intelligence and generation \citep{huang2023survey,zou2025fragment}, as well as focused accounts of layout generation \citep{shi2023layoutsurvey} and vector graphics \citep{tian2024vectorsurvey}.
Their task- and representation-oriented analyses provide essential context, but a cross-cutting comparison must also distinguish what a generated artifact makes possible from what an editing workflow has actually demonstrated or measured.
This survey addresses that distinction through a representation--editability--evaluation (\ac{rev}) framework for generative graphic design.
We examine delivered artifact types, the revision operations they support, and the evidence used to evaluate those operations, before identifying outstanding challenges in instruction adherence, reflow, preservation, artifact validity, and human workflows.
Figures~\ref{fig:rev} and \ref{fig:protocol} illustrate the framework and the associated evaluation questions.
